\documentstyle[%


righttag%


,11pt%


,amssymb%


,russian%               remove in Eng.


,izvru%                 Set izven% in Eng.


]{amsart}





%





\newcommand{\la}{\langle}


\newcommand{\ra}{\rangle}


\newcommand{\tts}[1]{}








%\hoffset=-15mm%                  For Eng.


\setcounter{page}{3}


\god{2005}


\nomer{11~(522)} %                        Remove in Eng.


%% remove ^^^^^^


%\nomer{Vol.\,49, No.\,11}%               Set in Eng.


%\str{\pageref{begin}--\pageref{end}}%  Set in Eng.


\udk{519.71}





\author{\it Н.З.\,Габбасов, Б.C.\,Кочкарев}


% NB:   ^^ remove in Eng.


\title{Реализация дискретных вероятностных


процессов конечными линейными и вероятностными автоматами}





\date{}


%\pagestyle{myheadings}%         Set in Eng.


\pagestyle{plain} %              Remove in Eng.


\begin{document}


%\thispagestyle{plain}%          Set in Eng.


\maketitle


\label{begin}








Пусть $X$ --- конечный непустой алфавит. Через $X^*$ обозначим


множество всех слов конечной длины над $X$, т.\,е. $X^*$ --- свободный 


моноид, генерируемый множеством $X$. Единицу моноида $X^*$


будем обозначать через $e$, т.\,е. $e$ --- пустое слово.


Если $p=x_1x_2\dotsc x_n$,


где $x_i\in X$, $i=\ov{1,n}$, то $|p|=n$ и


$p^\top=x_n x_{n-1}\dotsc x_2x_1$.





\begin{defn}


Всякое отображение $f:X^*\to R$, где $R$ --- поле,


будем называть дискретным процессом (кратко процессом).


\end{defn}





\begin{defn}


Сопряженным процессом к процессу $f$ будем называть 


процесс $f^\top$, определяемый условием


$$


f^\top(p)=f(p^\top)\ \ \forall p\in X^*.


$$


\end{defn}





\begin{defn}


Если $f=f^\top$, то процесс $f$ называется самосопряженным.


\end{defn}





\begin{note}


Очевидно, $(f^\top)^\top=f$.


\end{note}





\begin{defn}


Конечным линейным автоматом (КЛА) с выходом над $X$ будем


называть тройку


$$


A=\la\mu,\{A(x),\ x\in X\},\nu\ra,


$$


где $\mu$ --- $1\times n$-матрица над полем $R$, $A(x)$ ($x\in X$)


--- квадратные


$n\times n$-матрицы над полем $R$, $\nu$ --- $n\times1$-матрица над полем $R$.


\end{defn}





Очевидно, со всяким КЛА $A$ ассоциируется некоторый процесс


$f_A$, определяемый равенством


$$


f_A(p)=\mu A(p)\nu\ \ \forall p\in X^*,


$$


где $A(p)=A(x_1)A(x_2)\dotsc A(x_k)$, если $p=x_1x_2\dotsc x_k$,


$A(e)$ --- единичная $n\times n$-матрица.





В этом случае будем также говорить, что КЛА реализует процесс $f_A$.





\begin{defn}


Процесс $f$ будем называть линейным процессом


(ЛП), если существует КЛА $A$ такой, что $f_A=f$.


\end{defn}





Если $n=n_A$ --- минимальное возможное число состояний КЛА


$A$, реализующего ЛП $f$, то будем писать $n=\dim f$.





\begin{defn}


Процесс $f$ будем называть псевдовероятностным 


процессом (ПП), если $\sum\limits_{x\in X}f(px)=f(p)$ \; $\forall p\in X^*$.


\end{defn}





\begin{defn}


ПП $f$ будем называть вероятностным процессом 


(ВП), если поле $R$ вещественное и 1)~$f(e)=1$,


2)~$f(p)\ge0$ \ $\forall p\in X^*$.


\end{defn}





\begin{defn}


Если для ВП $f$ \ $\dim f=n<\infty$, то $f$ будем называть 


вероятностным линейным процессом (ВЛП).


\end{defn}





\begin{defn}


Конечным вероятностным автоматом (КВА) с


выходом над $X$ будем называть тройку


$$


B=\la\alpha,\{B(x),\ x\in X\},\beta\ra,


$$


где $\alpha$ --- стохастическая $1\times n$-матрица, $B(x)$ ($x\in B$)


--- квадратные


неотрицательные (поэлементно) $n\times n$-матрицы такие, что


$\sum\limits_{x\in X}B(x)$ есть стохастическая матрица, $\beta$ ---


$n\times1$-матрица, все


компоненты которой равны единице.


\end{defn}





Очевидно, со всяким КВА $B$ ассоциируется ВП


$$


f_B(p)=\alpha B(p)\beta\ \ \forall p\in X^*.


$$


В этом случае будем также говорить, что КВА $B$ реализует ВП $f_B$.





\begin{defn}


Если существует КВА $B$, который реализует


ВП $f$, то $f$ называется конечным вероятностным процессом


(КВП).


\end{defn}





Очевидно, всякий КВП $f$ является ВЛП, но обратное, как мы


увидим, не всегда имеет место.





Минимальное возможное число состояний КВА $B$, 


реализующего КВП $f$, обозначим через $\dim^+f$.





\begin{thm}


Существует ВП $f$ такой, что $\dim f=3$, но $\dim^+f=\infty$, $|X|=2$.


\end{thm}





\begin{pf}


Рассмотрим линейный автомат


$$


A=\la\mu,\{A(0),\ A(1)\},\nu\ra,


$$


где


\begin{gather*}


\mu=\frac{1}{\sqrt{1+\varepsilon^2}}(0,\varepsilon,1),\quad


A(0)=\frac{1}{2}


\begin{pmatrix}


\cos\theta & -\sin\theta & 0 \\


\sin\theta & \cos\theta & 0 \\


0 & 0 & 1


\end{pmatrix}


,\\


A(1)=\frac{1}{2(1+\varepsilon^2(2-\cos\theta))}


\begin{pmatrix}


\varepsilon^2\sin^2\theta &


\varepsilon^2(2-\cos\theta)\sin\theta & \varepsilon\sin\theta \\


\varepsilon^2(2-\cos\theta)\sin\theta &


\varepsilon^2(2-\cos\theta)^2 &


\varepsilon(2-\cos\theta)\\


\varepsilon\sin\theta & \varepsilon(2-\cos\theta) & 1


\end{pmatrix}


,\quad \nu=\frac{1}{\sqrt{1+\varepsilon^2}}


\begin{pmatrix}0 \\ \varepsilon \\ 1 \end{pmatrix}.


\end{gather*}


Непосредственными вычислениями показывается, что $f_A$ при


$|\varepsilon|\le\frac{1}{3}$ является ВП, причем в силу известного [1] критерия индуцируемости 


стохастического процесса марковской цепью, при условии, 


что $\varepsilon\ne0$ и $\theta$ --- угол иррационального порядка относительно $\pi$


(т.\,е. $\theta=r\pi$, где $r$ иррациональное), $f_A$ не может быть реализован 


никаким КВА $B$, т.\,е. в этом случае $\dim^+f=\infty$. Нетрудно


показать, что $\dim f_A=3$.


\end{pf}





\begin{thm}


Для любого $N\ge3$ существует ВП $f_N$ такой, что


$\dim^+f_N=N$, $\dim f_N=3$, $|X|=2$.


\end{thm}





\begin{pf}


Искомый ВП $f_N$ есть процесс, реализуемый


КЛА из теоремы~1 для $\theta$ подходящего рационального порядка 


относительно $\pi$.


\end{pf}





\begin{thm}


Всякий ВП $f$, для которого $\dim f<3$, реализуется


КВА, т.\,е. является КВП.


\end{thm}





\begin{pf}


В силу известного [1] критерия индуцируемости


стохастического процесса марковской цепью достаточно использовать 


тот факт, что острый выпуклый конус на двумерной плоскости 


имеет не более двух базисных образующих, а угол соответствующего 


сектора меньше~$\pi$.


\end{pf}





\begin{thm}


Процесс $f$ является ЛП тогда и только тогда, когда


$f^\top$ является ЛП, причем $\dim f=\dim f^\top$.


\end{thm}





\begin{pf}


{\bf Необходимость.} Пусть $f$ реализуется КЛА


$$


A=\la\mu,\{A(x),\ x\in X\},\nu\ra,


$$


т.\,е. $f(p)=\mu A(p)\nu$ \ $\forall p\in X^*$. Тогда


$f^\top(p)=f(p^\top)=\mu A(p^\top)\nu=\mu A(x_k)


\dotsc A(x_1)\nu$.


Очевидно, имеем


$$


f^\top(p)=f(p^\top)=\nu^\top A^\top(x_1)A^\top(x_2)\dotsc


A^\top (x_k)\mu^\top,


$$


где $A^\top(x)$ --- транспонированная матрица $A(x)$. Таким образом,


$f^\top$ реализуется посредством КЛА


$$


A^\top=\la\nu^\top,\{A^\top(x),\ x\in X\},\mu^\top\ra.


$$





{\bf Достаточность} в силу замечания 1 доказывается аналогично. Из


самого доказательства видно, что в этом случае $\dim f=\dim f^\top$.


\end{pf}





\begin{thm}


Если известно, что $\dim f\le n$, где $n$ задано и существует 


эффективный способ вычисления $f(p)$ для любого


$p$, $|p|\le2n-1$, то самосопряженность ЛП $f$ алгоритмически разрешима.


\end{thm}





\begin{pf}


Действительно, в этом случае достаточно проверить 


справедливость соотношения


$$


f(p)=f(p^\top)


$$


для всех слов длины $\le2n-1$.


\end{pf}





\begin{defn}


ПП $f$ называется двусторонним, если $f^\top$ является ПП.


\end{defn}





\begin{defn}


ВП $f$ называется двусторонним, если $f^\top$ является ВП.


\end{defn}





\begin{thm}


Для того чтобы ПП $f$ являлся двусторонним, необходимо 


и достаточно, чтобы выполнялось условие


$$


\sum_{x\in X}f(xp)=f(p)\ \ \forall p\in X^*.


$$


\end{thm}





\begin{pf}


{\bf Необходимость.} Пусть $f^\top$ является ПП. Тогда


$$


f(p)=f^\top(p^\top)=\sum_{x\in X}f^\top(p^\top x)=


\sum_{x\in X}f(xp)\ \ \forall p\in X^*.


$$





{\bf Достаточность.} Пусть теперь ПП $f$ удовлетворяет условию


$\sum\limits_{x\in X}f(xp)=f(p)$ \ $\forall p\in X^*$. Тогда имеем


$$


\sum_{x\in X}f^\top(px)=\sum_{x\in X}f(xp^\top)=f(p^\top)=f^\top(p),


$$


т.\,е. $f^\top$ является ПП.


\end{pf}





\begin{thm}


Если двусторонний ПП $f$ является ВП, то $f$ является двусторонним ВП.


\end{thm}





\begin{pf}


Поскольку ПП $f$ является ВП, то


$f^\top(e)=f(e^\top)=f(e)=1$, $f^\top(p)=f(p^\top)\ge0$ \ $\forall p\in X^*$. Так как $f$ является 


двусторонним ПП, то и $f^\top$ является ВП, т.\,е. $f$ является


двусторонним ВП.


\end{pf}





Очевидно, всякий самосопряженный ПП (ВП) $f$ является двусторонним 


ПП (ВП).





\begin{thm}


Существует самосопряженный ВП $f$ такой, что $\dim f=3$, но $\dim^+f=\infty$,


$|X|=2$.


\end{thm}





\begin{pf}


Рассмотрим линейный автомат


$$


A'=\la\mu,A(0),A'(1),\nu\ra


$$


где $\mu$, $\nu$, $A(0)$ такие же, что и при определении линейного автомата


из теоремы 1, а $A'(1)$ представляет собой матрицу вида


$$


A'(1)=\frac{1}{2R}


\begin{pmatrix}


\varepsilon^2\sin^2\theta &


\varepsilon^2(2-\cos\theta)\sin\theta & \varepsilon\sin\theta \\


-\varepsilon^2(2-\cos\theta)\sin\theta &


\varepsilon^2(2-\cos\theta)^2 &


\varepsilon(2-\cos\theta)\\


-\varepsilon\sin\theta & \varepsilon(2-\cos\theta) & 1


\end{pmatrix}


$$


для $R=\varepsilon^2(2-\cos\theta)+1$.





Непосредственными вычислениями и используя теорему 5 показываем,


что $f_{A'}$ при $|\varepsilon|\le\frac{1}{3}$ является самосопряженным ВП. По


той же причине, что и при доказательстве теоремы 1, $f_{A'}$ при условии, 


что $\varepsilon\ne0$ и $\theta$ --- угол иррационального порядка относительно


$\pi$, не может быть реализован никаким КВА $B$, т.\,е. в этом случае


$\dim^+f_{A'}=\infty$. Для завершения доказательства остается заметить,


что линейный автомат $A'$ является неприводимым, т.\,е. $\dim f_{A'}=3$.


\end{pf}





\begin{cor}


Существует двусторонний ВП $f$ такой, что


$\dim f=3$, но $\dim^+f=\infty$, $|X|=2$.


\end{cor}





\begin{thm}


Для любого $N\ge3$ существует самосопряженный ВП


$f_N$ такой, что\linebreak $\dim^+f_N=N$, $\dim f_N=3$, $|X|=2$.


\end{thm}





\begin{pf}


Искомый самосопряженный ВП $f_N$ есть процесс, 


реализуемый КЛА из теоремы 8 для $\theta$ подходящего рационального 


порядка относительно $\pi$.


\end{pf}





\begin{cor}


Для любого $N\ge3$ существует двусторонний ВП $f_N$


такой, что $\dim^+f_N=N$, $\dim f_N=3$, $|X|=2$.


\end{cor}





\begin{note}


В силу континуальности выбора значений параметра $\varepsilon$


для соответствующих ЛА в теоремах 1, 2, 8, 9 эти теоремы будут 


справедливы, если изменить их формулировки с учетом такого


выбора $\varepsilon$. Так, например, теорема 1 в этом случае примет 


следующий вид.


\end{note}





{\bf Теорема 1${}'$.} {\it Существует континуум ВП $f$ таких, что $\dim f=3$,


но $\dim^+f=\infty$, $|X|=2$.}





\begin{thebibliography}{9}





\bibitem{1}


Heller A. {\em On stochastic processes derived from Markov chains} // Ann.


Math. Statist. -- 1965. -- V.\,36. -- P.\,1286--1291.





\end{thebibliography}





\vskip 2cm





\post{\it Казанский государственный}{\it Поступила}


\post{\it университет}{26.09.2003}


\smallskip


\post{\it Татарский государственный}{}


\post{\it гумaнитарно-педагогический университет}{}





\label{end}


\end{document}


